\section{Three purely quadratic finite labels over the original base}
\label{sec:pure-quadratic-original-strict}

\paragraph{Scope.}
The three finite characters in this section have trivial cubic part.
Their strict groups are not consequences of the earlier nontrivial
cubic-sector arguments.  Six small biquadratic covers suffice.
The class input is an elementary ideal-count bound, and the unit
input consists of six explicitly checked Pell units.

Put $p=269$ and $K=\mathbb Q(\sqrt{-3})$.
For $d\in\{-4,8,-8\}$, write $\delta_d=\chi_d|_{G_K}$ for
the primitive rational quadratic character of discriminant $d$,
restricted to $K$.
Let $c_p$ be the actual cyclotomic character and $\omega_p$ its
Teichmuller part.  For a discrete module $D$ use
\[
 \begin{gathered}
 S_{\rm str}(K,D)=\ker\bigl(H^1_{\rm cts}(G_K,D)
                                      \longrightarrow\mathcal L(D)\bigr),\\
 \mathcal L(D)=\prod_{v\mid p}H^1(G_{K_v},D)
              \times\prod_{v\nmid p,\infty}H^1(I_v,D).
 \end{gathered}
 \tag{PQ.1}
\]

\begin{theorem}[AC-G160: the three pure quadratic original-coefficient rows]
\label{AC-G160}
Let $E/\mathbb Q_p$ be any finite extension with its original ring
of integers $O$.  Let $d\in\{-4,8,-8\}$, $m\in\mathbb Z$,
$e\geq1$, and let $\kappa:G_K\to1+\mathfrak m_O$ be continuous
and unramified at every finite place away from $p$.
If $e+m\equiv133$ or $267\pmod{268}$, then
\[
 \begin{gathered}
 D=(E/O)(\delta_d\omega_p^{-m}\kappa^{-1}c_p^{-e}),\\
                         S_{\rm str}(K,D)=0.
 \end{gathered}
 \tag{PQ.2}
\]
The residual strict group is zero over every finite field of
characteristic $p$.  The coefficient ring may have arbitrary
absolute ramification.
For $m=0,\kappa=1$, all three finite labels already occur on
the prescribed AC-K25 hosts.  They carry no nontrivial cubic factor.
\end{theorem}

\begin{proof}
\proofstep{1. Six exact residual covers.}
Let $k=O/\mathfrak m_O$, and let $a$ denote the residual character
of the coefficient in (PQ.2).
For $r=e+m\pmod{268}$, the quotient by the Kummer twist is
\[
 \eta=a/\bar c_p=
 \begin{cases}
   \delta_d,&r=267,\\
   \delta_d\bar\tau_p,&r=133,
 \end{cases}
 \qquad \tau_p=\omega_p^{134}.
 \tag{PQ.3}
\]
The quadratic Gauss identity identifies $\tau_p$ with the
character of $\mathbb Q(\sqrt p)$, since $p\equiv1\pmod4$;
this is the same exact identification used in AC-G154, Step1.
Put $\nu=0$ for $r=267$ and $\nu=1$ for $r=133$.
The complete field table is
\[
 \begin{array}{c|r|r|r}
 d&r_0&D_0&N\\ \hline
 -4&-1&3&K(\sqrt{-p^\nu})\\
  8& 2&2&K(\sqrt{2p^\nu})\\
 -8&-2&6&K(\sqrt{-2p^\nu}).
 \end{array}
 \tag{PQ.4}
\]
In each row $N=K(\sqrt{r_0p^\nu})$ is biquadratic over
$\mathbb Q$ and is exactly the quadratic cover of $K$ splitting
$\eta$.  None of $r_0p^\nu$ is $-3$ times a rational square.
Its real quadratic subfield is $\mathbb Q(\sqrt D)$, where
$D=D_0p^\nu$.
For a negative $r_0$ this follows by multiplying its square root
by $\sqrt{-3}$; for $r_0=2$ the field is already real.
Let $b$ generate $\operatorname{Gal}(N/K)$.
Then $b(\sqrt D)=-\sqrt D$ and $\eta^{-1}(b)=-1$.
The factor $\tau_p$ describes a residual cover only.
It is not added to the original finite label $\delta_d$.

\proofstep{2. A uniform elementary class bound.}
The three quadratic subfields of the six covers have the following
fundamental discriminants:
\[
 \begin{array}{c|c|r|r|r}
 d&\nu&\multicolumn{3}{c}{\text{quadratic discriminants}}\\ \hline
 -4&0&-3&-4&12\\
 -4&1&-3&-1076&3228\\
  8&0&-3&8&-24\\
  8&1&-3&2152&-6456\\
 -8&0&-3&-8&24\\
 -8&1&-3&-2152&6456.
 \end{array}
 \tag{PQ.5}
\]
These follow from the integral quadratic basis and the squarefree
radicands $-3,r_0p^\nu,-3r_0p^\nu$.
Every absolute discriminant is at most $6456$.
The Minkowski representative bound for a real quadratic field is
$\sqrt{|\Delta|}/2$, and for an imaginary quadratic field it is
$2\sqrt{|\Delta|}/\pi$.
Both are strictly smaller than $54$, because
\[
 \frac{2\sqrt{6456}}3<54,
 \qquad 4\cdot6456=25824<26244=9\cdot54^2,
 \qquad \pi>3.
 \tag{PQ.6}
\]
The real bound is smaller still, since $1/2<2/3$.
Thus each class has an integral ideal representative of norm
at most $53$.
The number of ideals of norm $n$ in a quadratic field is at most
$d(n)$: the split-prime coefficient is $a+1$ at a prime power
$q^a$, and inert or ramified primes give no larger coefficient.
Multiplication over primes gives the bound for every $n$.
Consequently all eighteen entries in (PQ.5) have
\[
 h\leq\sum_{n=1}^{53}d(n)
      =\sum_{a=1}^{53}\left\lfloor\frac{53}{a}\right\rfloor
      =219<269.
 \tag{PQ.7}
\]
No exact quadratic class number is needed.

For any nonidentity $h\in\operatorname{Gal}(N/\mathbb Q)=C_2^2$,
norm to its quadratic fixed field followed by extension is
$1+h$ on $\operatorname{Cl}(N)[p]$.
Its target class $p$-torsion is zero by (PQ.7), hence $h$ acts
by minus one.  The product of two distinct nonidentity elements
is the third; it would act both by plus one and minus one.
Since $p$ is odd, this proves
\[
                            \operatorname{Cl}(N)[p]=0
 \tag{PQ.8}
\]
for every cover in (PQ.4).
This is the complete elementary V4 argument, applied directly to
these fields rather than inferred from a cubic sector.

\proofstep{3. Full-local Kummer descent and its one unit line.}
The degree $[N:K]=2$ is prime to $p$.
Prime-to-$p$ restriction/inflation, also on each local decomposition
or inertia group, identifies the residual strict group with
\[
 \begin{gathered}
 e_-\ker\left(V_{N,S}\otimes_{\mathbb F_p}k
                   \longrightarrow\mathcal P_N\otimes_{\mathbb F_p}k\right),\\
 V_{N,S}=\{x\in N^\times:p\mid v(x)\text{ for }v\nmid p\}/N^{\times p},\\
 \mathcal P_N=\bigoplus_{v\mid p}N_v^\times/N_v^{\times p},
                         \qquad e_-=(1-b)/2.
 \end{gathered}
 \tag{PQ.9}
\]
Indeed the coefficient over $N$ is $k(1)$, and the extra quotient
twist selects $\eta^{-1}$, the minus character of $b$.
There is no need to adjoin $\mu_p$.
The full local power condition also forces the valuations above
$p$ to be divisible by $p$.
Thus its kernel lies in the all-primes valuation group $V_N$,
with the exact sequence
\[
 0\longrightarrow\mathcal O_N^\times/\mathcal O_N^{\times p}
 \longrightarrow V_N\longrightarrow\operatorname{Cl}(N)[p]
 \longrightarrow0.
 \tag{PQ.10}
\]
This is the same full-local Kummer sequence used in AC-G154,
Step3, and AC-G156, Step2.
By (PQ.8) only the global units remain.

The field $N$ has signature $(0,2)$, so its unit rank is one.
It has no root of unity of order $p$, since $268$ does not divide
$[N:\mathbb Q]=4$.
Its unit quotient modulo $p$ is therefore a one-dimensional
$\mathbb F_p$-space.
The norm to $K$ takes units into $\mathcal O_K^\times=\mu_6$,
which is zero in this quotient since $p\nmid6$.
Hence $1+b=0$ on the whole unit quotient: it is exactly the
minus line needed in (PQ.9).
It is enough to exhibit a single global unit with nonzero
local Kummer image for each cover.

\proofstep{4. Six explicit global units.}
Write $x=\sqrt D$ in the real quadratic subfield of $N$.
The units are $U=a+cx$, with data
\[
 \begin{array}{r|r|r|r}
 D&a&c&a^2-Dc^2\\ \hline
 3&2&1&1\\
 2&1&1&-1\\
 6&5&2&1\\
 807&51841948&1824923&1\\
 538&69051&2977&-1\\
 1614&291698879051&7260775270&1.
 \end{array}
 \tag{PQ.11}
\]
Every row is checked by integer multiplication in the supplied
finite certificate.  If $\epsilon=a^2-Dc^2\in\{1,-1\}$,
then $U$ satisfies $X^2-2aX+\epsilon=0$.
Thus $U$ is integral, and $U^{-1}=(a-cx)/\epsilon$ is integral.
No assertion that $U$ is a fundamental unit is used.
The actual global quotient
\[
                     V=U/b(U)=U^2/\epsilon
 \tag{PQ.12}
\]
is a unit in the minus component modulo $p$.
All local tests below apply to this global quotient before
localization, including when $b$ exchanges completions.

\proofstep{5. The three unramified tests.}
For $\nu=0$, every completion of $N$ above $p$ is unramified
of absolute degree two.  Indeed $K_p/\mathbb Q_p$ is the
unramified quadratic extension; adjoining a square root of a
rational $p$-adic unit adds no further extension, since its
quadratic extension is either split or this same unramified one.
Thus $N$ has two $p$-places, each of ramification index one.
The automorphism $b$ may exchange them; no action on an isolated
completion is assumed.

Let $q=p^2-1=72360$.
Exact repeated squaring in
$(\mathbb Z/p^2\mathbb Z)[X]/(X^2-D)$ gives
\[
 \begin{array}{c|r|c}
 D&B&U^q\pmod{p^2}\\ \hline
 3&158&1+158p x\\
 2&18&1+18p x\\
 6&13&1+13p x.
 \end{array}
 \tag{PQ.13}
\]
The displayed integers are finite ring computations, retaining
both coefficients throughout.
Since $\epsilon^q=1$, (PQ.12) gives
$V^q\equiv1+2Bp x\pmod{p^2}$.
At any completion, $x$ is a unit and $B\not\equiv0\pmod p$.
Therefore $V^q-1$ has valuation exactly one.

In an unramified $p$-adic field, logarithm and exponential are
inverse between principal units and the maximal ideal.
The logarithm of $V^q$ also has valuation exactly one: the
terms of degree at least two have strictly higher valuation.
A $p$th power of a principal unit has logarithm of valuation
at least two.  Moreover if a principal unit has a $p$th root
among all local units, that root has residue one because
Frobenius on the finite residue field is bijective.
If $V$ were a local $p$th power, $V^q$ would be such a principal
$p$th power, contradicting its logarithm valuation.
Thus (PQ.13) supplies nonzero full local Kummer images of the
three actual global units $V$.

\proofstep{6. The three ramified tests.}
For $\nu=1$, $D$ is $3p,2p$ or $6p$.
The element $x=\sqrt D$ is a uniformizer of the ramified quadratic
extension of $K_p$, with $v_x(p)=2$.
This completion has absolute degree four, equal to $[N:\mathbb Q]$,
so $N$ has a unique $p$-place, with indices $(2,2)$.
Every constant $a$ in the final three rows of (PQ.11) is a
$p$-adic unit.  Direct reduction gives
\[
 \begin{array}{c|r|r}
 D&c/a\pmod p&2c/a\pmod p\\ \hline
 807&242&215\\
 538&131&262\\
 1614&169&69.
 \end{array}
 \tag{PQ.14}
\]
The quotient $V=(a+cx)/(a-cx)$ is a principal unit, and
\[
                     \log V=(2c/a)x+O(x^2).
 \tag{PQ.15}
\]
Each leading coefficient in (PQ.14) is nonzero.
Because the absolute ramification index is $2<p-1$, logarithm
and exponential again identify all principal units with the
maximal ideal.  A $p$th power has logarithm valuation at least
$2+1=3$, whereas (PQ.15) has valuation one.
A root of a principal unit has residue one as in Step5.
Consequently all three classes are nonzero in the full local
power quotient, not merely in a quotient by a smaller subgroup.

\proofstep{7. Residual vanishing and the original coefficient lift.}
Each displayed global unit $V$ has nonzero local Kummer image.
Since the unit quotient is the one-dimensional minus line,
localization is injective on that line; it remains so after
extending scalars to $k$.
Equations (PQ.8)--(PQ.10) prove the residual strict vanishing.

Let $\varpi$ be a uniformizer of the original $O$.
Return to the coefficient $D$ of (PQ.2), and let $a$ again denote
its residual character from Step1.  Then $D$ is divisible and
$D[\varpi]\simeq k(a)$.
At $p$, the finite character $\delta_d$ is unramified.
The residual inertia action is $\bar c_p^{-267}$ or
$\bar c_p^{-133}$; both have exact order $268$, since both
exponents are coprime to $268$ and $K_p/\mathbb Q_p$ is unramified.
The character $\kappa$ reduces to one.
Some inertia element therefore acts by a scalar whose difference
from one is a unit of $O$, so $D^{G_{K_p}}=0$.
At every finite place away from $p$, the cyclotomic, Teichmuller
and $\kappa$ factors are unramified, and the remaining inertia
character has order dividing two.
Its invariants are zero or the whole divisible module $D$.
Thus $D^J/\varpi D^J=0$ for every local group $J$ in (PQ.1).

The multiplication-by-$\varpi$ long exact sequence gives local
injectivity $H^1(J,D[\varpi])\to H^1(J,D)$.
Every global strict class killed by $\varpi$ lifts to a residual
class, and local injectivity makes the lift strict.
It is zero by the residual assertion just proved.
Finally every continuous cocycle from compact $G_K$ to discrete
$E/O$ has finite image, killed by some power of $\varpi$.
Any nonzero strict class would produce a nonzero one killed by
$\varpi$, a contradiction.  This proves (PQ.2) for the original,
possibly ramified coefficient ring.
\end{proof}

\paragraph{Actual finite labels and exact count.}
For $m=0,\kappa=1$ the three finite labels are precisely
$\delta_{-4},\delta_8,\delta_{-8}$.
Their fields are contained in $\mathbb Q(\zeta_8)$ and hence
in every prescribed host with $u\geq2$.
Their restrictions to $K$ are distinct and nontrivial, because
$K\cap\mathbb Q(\zeta_8)=\mathbb Q$.
Their nontriviality separates them from the trivial finite label.
Restriction to the cubic subgroup separates them from $\chi$,
$\chi^2$ and every other label with a nontrivial cubic factor.
The positive weights are exactly $e=134n-1$, $n\geq1$;
there are
\[
                             3\lfloor T/134\rfloor
 \tag{PQ.16}
\]
pairs of a displayed finite label and $1\leq e<T$.
Only full strict vanishing is claimed here.
The positive derived indices retain their local comparison cost,
and the common-window implication requires its separate proof.

\paragraph{Evidence and source scope.}
The self-contained finite calculation is
\path{verify_quadratic_units.py}, with output
\path{quadratic_units_certificate.json}.
It verifies the six exact unit norms, all discriminants in (PQ.5),
the bound (PQ.7), all three modular powers and all three ramified
logarithm coefficients.
The Pell search merely finds the displayed integers; every proof
input is rechecked by finite arithmetic.
There is no analytic interval or class-number selector.
Minkowski and class extension/norm are used in precisely the
forms already audited in AC-G154, Step2.
The full-local Kummer sequence retains Stacks tag~03PK and the
prime-to-$p$ descent used in AC-G154, Step3, and AC-G156, Step2.
Dirichlet's unit rank formula is used only for this biquadratic
rank-one case.
Geometrically the coefficients are the original rank-one
Artin--Tate systems on the arithmetic curve of $K$, with finite
covers as in SGA~1, Expos\'e~V, \S\S4 and7, at the previously
audited scope.  No new Grothendieck archive attribution, mixed
realization, all-prime index bound, H04.4 or cubic RH follows.
No Lean execution is used.
